% ID: SEL-D0018-BC-E06
% ESTADO: cotejada. Segunda pasada visual completa y borrador compilado revisados.
% Fuente: Selectividad/matematicas_aplicadas/matematicas_aplicadas_P_F.pdf
% SHA-256: fe6fd5b48407569d41169c80d71e6997050d6d06d1afd41cd639a553f9b712b5
% Página: 2; bloque C; ejercicio original 6.

El 47\% de los jóvenes andaluces tienen una vida sedentaria. De ellos, el 72\% presentan obesidad, mientras que solamente la presentan el 22\% de los jóvenes no sedentarios. Se elige al azar un joven andaluz.
\begin{enumerate}[label=\alph*)]
\item \textbf{(1 punto)} Calcule la probabilidad de que sea sedentario y no presente obesidad.
\item \textbf{(0.75 puntos)} Calcule la probabilidad de que presente obesidad.
\item \textbf{(0.75 puntos)} Calcule la probabilidad de que sea sedentario, sabiendo que presenta obesidad.
\end{enumerate}
